\documentclass[11pt,a4paper]{article}

% ===== Paquetes básicos =====
\usepackage[spanish,es-noquoting]{babel} % español
\usepackage[utf8]{inputenc}              % UTF-8
\usepackage[T1]{fontenc}                 % acentos correctos en PDF
\usepackage{lmodern}                     % mejor tipografía
\usepackage{geometry}                    % márgenes
\geometry{margin=2.5cm}
\usepackage{amsmath,amssymb,amsthm}      % mates
\usepackage{hyperref}                    % links clickeables
\hypersetup{colorlinks=true, linkcolor=blue, urlcolor=blue}
\usepackage{enumitem}                    % listas compactas
\setlist{itemsep=0.2em, topsep=0.2em}

% ===== Comandos útiles =====
\newcommand{\N}{\mathbb{N}}
\newcommand{\Pof}[1]{\mathcal{P}(#1)}    % conjunto de partes
\newcommand{\compl}[1]{\overline{#1}}    % complemento

\title{Apuntes de Álgebra de Matemática Discreta}


\begin{document}
\maketitle
\tableofcontents
\bigskip
\newpage

\section{Unidad 3: Álgebras de Boole}

\subsection{Conjunto de partes}
Dado un conjunto $S$, el \textit{conjunto de partes} $\Pof{S}$ es el conjunto de todos los subconjuntos de $S$.

Sea $S=\{a,b\}$. Entonces:
\[
\Pof{S}=\{\varnothing,S\}
\quad\text{(conjunto trivial)}
\]
Más explícitamente:
\[
\Pof{S}=\{\varnothing,\{a\},\{b\},S\}.
\]

\textbf{Nota.} Todos los elementos de $\Pof{S}$ se escriben entre llaves porque son conjuntos, excepto el conjunto vacío $\varnothing$.

\subsection{Álgebra de Boole de conjuntos}

\subsubsection{Relación de divisibilidad}
Si $x,y\in\N$, se dice que “$x$ divide a $y$” si y sólo si:
\[
x \mid y \;\Longleftrightarrow\; \exists\,k\in\N:\; y=kx.
\]

\subsubsection{Número natural libre de cuadrados}
Un número $n\in\N$ es \textit{libre de cuadrados} si ningún cuadrado mayor que $1$ lo divide.

\subsection{Álgebra de Boole de los divisores de n}
Sea un número $n\in\N$ libre de cuadrados, entonces $(D(n), \lor, \land, \lnot, 1, n)$ es un \textit{Álgebra de Boole},
donde la relación enttre sus elementos es la \textit{\textbf{divisibilidad}}.

\begin{itemize}
    \item $D(n)$ es el conjunto de todos los divisores de $n$.
    \item $x \lor y = \text{mcm}(x,y)$.
    \item $x \land y = \text{mcd}(x,y)$.
    \item $\lnot x = \frac{n}{x}$.
    \item $1$ es el elemento neutro para $\lor$.
    \item $n$ es el elemento neutro para $\land$.
\end{itemize}

\subsection{Propiedades}
\paragraph{Doble complemento.}
\[
\compl{\compl{A}}=A, \qquad
\lnot(\lnot)x = x
\]

\paragraph{Leyes de De Morgan.}
\[
\compl{A\cup B} = \compl{A}\cap\compl{B},\qquad
\lnot(x\lor y) = \lnot x \land \lnot y
\]

\[
\compl{A\cap B} = \compl{A}\cup\compl{B},\qquad
\lnot(x\land y) = \lnot x \lor \lnot y
\]

\paragraph{Leyes conmutativas}
\[
A \cup B = B \cup A,\qquad
x \lor y  = y \lor x
\]

\[
A \cap B = B \cap A, \qquad
x \land y = y \land x
\]

\paragraph{Leyes asociativas}
\[
A \cup (B \cup C) = (A \cup B) \cup C,\qquad
x \lor (y \lor z) = (x \lor y) \lor z
\]

\[
A \cap (B \cap C) = (A \cap B) \cap C, \qquad
x \land (y \land z) = (x \land y) \land z
\]

\paragraph{Leyes distributivas}
\[
A \cup (B \cap C) = (A \cup B) \cap (A \cup C),\qquad
x \lor (y \land z) = (x \lor y) \land (x \lor z)
\]

\[
A \cap (B \cup C) = (A \cap B) \cup (A \cap C), \qquad
x \land (y \lor z) = (x \land y) \lor (x \land z)
\]

\paragraph{Leyes de idempotencia}
\[
A \cup A = A,\qquad
x \lor x = x
\]

\[
A \cap A = A, \qquad
x \land x = x
\]

\paragraph{Leyes de identidad}
\[
A \cup \varnothing = A,\qquad
x \lor 1 = x
\]

\[
A \cap S = A, \qquad
x \land n = x
\]

\paragraph{Leyes de los inversos}
\[
A \cup \compl{A} = S,\qquad
x \lor \lnot x = n
\]

\[
A \cap \compl{A} = \varnothing, \qquad
x \land \lnot x = 1
\]

\paragraph{Leyes de dominación}
\[
A \cup S = S,\qquad
x \lor n = n
\]

\[
A \cap \varnothing = \varnothing, \qquad
x \land 1 = 1
\]

\paragraph{Leyes de absorción}
\[
A \cup (A \cap B) = A,\qquad
x \lor (x \land y) = x
\]

\[
A \cap (A \cup B) = A, \qquad
x \land (x \lor y) = x
\]

\subsection{Álgebra de los interruptores}
\paragraph{Circuito en serie.} $x \cdot y = xy$
\paragraph{Circuito en paralelo.} $x + y = xy$

\subsection{Funciones de conmutación}
Sea $B = {0, 1}$, el conjunto de los valores booleanos.

\paragraph{Suma} $$x + y = \begin{cases}
0 & \Leftrightarrow x = 0 \boldsymbol{\land} y = 0 \\
1 & \Leftrightarrow x = 1 \lor y = 1
\end{cases}$$

\paragraph{Producto} $$x \cdot y = \begin{cases}
0 & \Leftrightarrow x = 0 \lor y = 0 \\
1 & \Leftrightarrow x = 1 \boldsymbol{\land} y = 1
\end{cases}$$

\paragraph{Negación} $$\overline{x} = \begin{cases}
0 & \Leftrightarrow x = 1 \\
1 & \Leftrightarrow x = 0
\end{cases}$$

\subsection{Funciones booleanas}

Sea $n \in \mathbb{Z}^+$, $B^n = {(b_1, b_2, ..., b_n) / b_i \in\ {0, 1}, \forall i}$
($B_n$ se forma de n-uplas de 0 o 1).
Una función $f: B^n \rightarrow B$ se llama función booleana o \textit{\textbf{función de conmutación}}
y su tabla contiene $2^n$ filas.

\end{document}
